% !TeX root = ../manual.tex
% ===========================================================================
%  Guía para el usuario final
% ===========================================================================

\section{Guía para el usuario final}
\label{sec:usuario-final}

Al iniciar sesión como usuario final, el menú lateral muestra dos
secciones: \boton{Mis Reservas} y \boton{Disponibilidad}.

\subsection{Consultar disponibilidad}
\label{sec:usuario:disponibilidad}

Esta es la pantalla desde la que se reserva. Muestra, para un deporte y una
fecha, todas las canchas de ese deporte y sus bloques horarios.

\begin{enumerate}
  \item Entra a \boton{Disponibilidad} desde el menú lateral.
  \item Elige el deporte: \boton{Pádel}, \boton{Tenis} o \boton{Básquet}.
  \item Elige el día, entre hoy y los cuatro días siguientes.
\end{enumerate}

\captura{usuario-disponibilidad}{1}{Disponibilidad de las canchas de un deporte, para el día seleccionado.}{usuario-disponibilidad}

Cada cancha aparece con su nombre, su horario de atención y una fila de
bloques horarios de una hora. El color del bloque indica su estado, según
la leyenda de la esquina superior derecha:

\begin{itemize}
  \item \textbf{Libre} (verde): se puede reservar.
  \item \textbf{Ocupado} (gris): ya tiene una reserva confirmada.
  \item \textbf{Mantenimiento} (naranja): el administrador bloqueó ese
        rango; no admite reservas nuevas aunque esté vacío.
\end{itemize}

\subsection{Crear una reserva}
\label{sec:usuario:reservar}

\begin{enumerate}
  \item En \boton{Disponibilidad}, pulsa un bloque \textbf{libre} de la
        cancha que quieras. El bloque se marca en azul y aparece el botón
        \boton{Reservar HH:MM} junto al nombre de la cancha.
\end{enumerate}

\captura{usuario-seleccion-bloque}{1}{Bloque libre seleccionado: aparece el botón para reservarlo.}{usuario-seleccion-bloque}

\begin{enumerate}
  \setcounter{enumi}{1}
  \item Pulsa \boton{Reservar HH:MM}. El sistema muestra una pantalla de
        confirmación con la cancha, la fecha, el horario y a nombre de
        quién queda la reserva.
\end{enumerate}

\captura{usuario-confirmar-reserva}{1}{Confirmación de los datos antes de crear la reserva.}{usuario-confirmar-reserva}

\begin{enumerate}
  \setcounter{enumi}{2}
  \item Revisa los datos y pulsa \boton{Confirmar Reserva}. Si prefieres no
        continuar, pulsa \boton{Cancelar} para volver a Disponibilidad sin
        reservar nada.
\end{enumerate}

Al confirmar, aparece un aviso de \textit{«Reserva Confirmada»} y el
sistema te lleva a \boton{Mis Reservas}, donde la nueva reserva figura
como \textsc{Confirmada}.

\aviso{Solo puedes tener confirmadas, al mismo tiempo, hasta el
\textbf{tope de reservas activas} que el club tenga configurado (3 de
forma predeterminada). El administrador puede cambiar ese número desde
\cref{sec:admin:configuracion}; si ya llegaste al tope vigente, cancela
alguna reserva antes de crear una nueva.}

\aviso{Si otra persona reserva el mismo bloque un instante antes que tú, tu
solicitud se rechaza con un mensaje en pantalla y el bloque pasa a verse
como ocupado. Es normal: el sistema garantiza que un mismo bloque nunca
quede reservado dos veces.}

\subsection{Ver Mis Reservas}
\label{sec:usuario:mis-reservas}

\boton{Mis Reservas} lista todas tus reservas, pasadas y futuras. Los
filtros de la parte superior —\boton{Todas}, \boton{Confirmadas},
\boton{Finalizadas}, \boton{Canceladas}— acotan la lista por estado:

\begin{itemize}
  \item \textsc{Confirmada}: la reserva es válida y todavía no ha
        comenzado su bloque horario. Se puede cancelar.
  \item \textsc{Finalizada}: el bloque horario de la reserva ya pasó. No
        se puede cancelar.
  \item \textsc{Cancelada}: la reserva se canceló, por ti o por el
        administrador. El bloque quedó libre nuevamente.
\end{itemize}

\captura{usuario-reservas-finalizadas}{1}{Filtro «Finalizadas»: reservas cuyo bloque horario ya pasó.}{usuario-reservas-finalizadas}

\subsection{Cancelar una reserva}
\label{sec:usuario:cancelar}

\begin{enumerate}
  \item En \boton{Mis Reservas}, localiza la reserva \textsc{Confirmada}
        que quieres cancelar.
  \item Pulsa \boton{Cancelar} en esa tarjeta.
\end{enumerate}

La reserva cambia de inmediato a estado \textsc{Cancelada} y el bloque
horario vuelve a aparecer como libre en \boton{Disponibilidad}.

\captura{usuario-reserva-confirmada}{1}{Antes de cancelar: la reserva está \textsc{Confirmada} y muestra el botón Cancelar.}{usuario-reserva-confirmada}

\captura{usuario-reserva-cancelada}{1}{Después de cancelar: la reserva pasa a \textsc{Cancelada}.}{usuario-reserva-cancelada}

\aviso{Una reserva cuyo bloque horario ya pasó (estado \textsc{Finalizada})
no se puede cancelar, aunque en su momento haya estado confirmada.}
